\newpage
\begin{center}{\heiti\large 附录 A\quad 补充材料与程序}\end{center}
\addcontentsline{toc}{section}{附录 A 补充材料与程序}
\setcounter{subsection}{0}
\renewcommand{\thesubsection}{A.\arabic{subsection}}

\subsection{支撑材料清单与运行说明}

正文各结论均按证据分层陈述：解析推导、拟合估计、尺度标定、留出检验、独立检验、评分后诊断、半合成数据、外推或估算，以及历史口径对照，其性质与边界已在正文相应处注明；半合成、估算与外推数据不作为观测验证。正文中 \texttt{*.csv}、\texttt{*.json}、\texttt{*.npz} 形式的文件名是相应结果的落盘名称：数据表在 \texttt{src/outputs\_qN/tables/}，跨问接口在 \texttt{src/outputs\_qN/interface/}，图在 \texttt{src/outputs\_qN/figures/}；每张表与每幅图的标题内注明其数据来源文件。

问题一：\texttt{src/code\_q1/q1\_00\_common.py}（公共路径与工具）$\to$ \texttt{q1\_01\_extract.py} $\to$ \texttt{q1\_02\_quality.py} $\to$ \texttt{q1\_03\_conflict.py}、\texttt{q1\_04\_bridge.py} $\to$ \texttt{q1\_05\_mixture.py} $\to$ \texttt{q1\_06\_scale.py}、\texttt{q1\_07\_optimal.py} $\to$ \texttt{q1\_08\_audit.py}；环境 Python 3.12.13、NumPy 2.4.4、pandas 3.0.2，不依赖 SciPy 与 scikit-learn。

问题二：\texttt{src/code\_q2/q2\_01\_fit.py}、\texttt{q2\_02\_theory.py}、\texttt{q2\_03\_economics.py}；环境 Python 3.12.7、NumPy 1.26.4、pandas 2.2.2、SciPy 1.13.1，随机种子 $2026$；只读 \texttt{src/outputs\_q1/interface/}，不读附件 A 原始表。

问题三：\texttt{src/code\_q3/q3\_00\_model.py}、\texttt{q3\_01\_optimize.py}、\texttt{q3\_02\_transition.py}、\texttt{q3\_03\_mixture\_joint.py}；已按幂次目标式与固定配比乘子 $m_*\approx0.8989423$ 重跑，P3 主网格为 210 行（运行清单 \texttt{P3\_run\_manifest.json}，2026-09-26 11:37）；运行清单的输入/输出哈希因后续上游文件更新需重新生成。

问题四：\texttt{src/code\_q4/q4\_01\_c8\_parse.py} $\to$ \texttt{q4\_02\_saturation.py} $\to$ \texttt{q4\_03\_bridge.py} $\to$ \texttt{q4\_04\_decomp.py} $\to$ \texttt{q4\_05\_forecast.py} $\to$ \texttt{q4\_99\_check\_report.py}；该自检（16 项）当前 3 项未通过：P3 输入哈希、P3 输出哈希一致性，以及 P4 目标日期（历史锚点 2026.25 已过去）。

注：以上脚本路径与随机种子以各问 Spec 的「当前实现差异」与修订历史为准；本节只登记，不复述 Spec。

\subsection{核心程序}

\textbf{(1) CRITIC 权重与加权 Huber 主分（\texttt{code\_q1/q1\_02\_quality.py} 节选）}
\begin{lstlisting}
def fit_critic(self, Z):
    sd = Z.std(axis=0, ddof=1)                 # contrast intensity
    R = np.corrcoef(Z.T)                       # indicator correlation
    C = sd * (1 - R).sum(axis=1)               # information content C_j
    self.critic_w = C / C.sum()
    return self.critic_w

def fit_delta(self, Z):                       # A1 weighted-MAD rule
    center = weighted_median(Z, self.critic_w)
    dev = np.abs(Z - center[:, None])
    mad = weighted_median(dev, self.critic_w)
    self.delta = float(np.clip(1.345 * 1.4826 * np.median(mad), .01, 1.0))
    return self.delta

def score(self, Z):                           # fixed weights and delta
    return weighted_huber_center(Z, self.critic_w, self.delta)
\end{lstlisting}

\textbf{(2) 问题三向量化求解器（\texttt{code\_q3/q3\_00\_model.py} 节选）}

\textbf{注：}下列节选即现行实现——幂次质量乘子 $Q^{-\kappa}$ 与固定配比乘子 $m_*=\exp\{s_{\rm red}\tilde\varphi(\bp^*)\}\approx0.8989423$，质量地板 $k_0(1-Q)$ 留在括号外；内层 $\log_{10}N$ 黄金分割与三级网格加密的结构不变，此处列出质量式与 $\Phi$ 判据函数。
\begin{lstlisting}
def hN(Q, p=PAR):                       # v5 power multiplier (drop linear 1+k(1-Q))
    return Q ** -p["kappa_N"]

def hD(Q, p=PAR):
    return Q ** -p["kappa_D"]

def loss(N, D, Q, p=PAR, m=M_STAR):
    """v5: floor k0(1-Q) outside bracket; mixture multiplier m only on scale terms"""
    return (p["E"] + p["k0"] * (1.0 - Q)
            + m * (p["A"] * hN(Q, p) * N ** -p["al"]
                   + p["B"] * hD(Q, p) * D ** -p["be"]))

def phi_ratio(C, g, Lctx, Qeval, Q0, p=PAR, eta=ETA, m=M_STAR):
    """v5 complete Phi at the fixed-Q optimal (N,D):
       Phi = [k0 + (m/Q)(kN TN + kD TD)] * S / [m * beta * TD * g'(Q)],
       S = (6+eta*L_ctx)N + dg,  TN = A N^-al Q^-kN,  TD = B D^-be Q^-kD"""
    kappa = 6 + eta * Lctx
    dg = max(g(Qeval) - g(Q0), 0.0)
    _, x = _profile(C, Qeval, dg, kappa, p, m=m)   # inner golden section (unchanged)
    N = 10.0 ** x; D = C / (kappa * N + dg)
    TN = p["A"] * N ** -p["al"] * hN(Qeval, p)       # TN = A N^-al Q^-kN
    TD = p["B"] * D ** -p["be"] * hD(Qeval, p)       # TD = B D^-be Q^-kD
    S = kappa * N + dg
    num = (p["k0"] + m / Qeval * (p["kappa_N"] * TN + p["kappa_D"] * TD)) * S
    den = m * p["be"] * TD * g_prime(Qeval)
    return float(num / den)             # C_crit solves Phi(C) = 1
\end{lstlisting}

\textbf{(3) 结构分解模型与 Shapley 分解（\texttt{code\_q4/q4\_04\_decomp.py} 节选）}
\begin{lstlisting}
def f_add(p, L, t):                            # M1b: S-bridge + linear tech term
    return p["lo"] + p["amp"] * sig(L, p["k"], p["L0"]) + p["tau"] * t

def shapley(fun, L0, t0, L1, t1):
    """two-factor Shapley: scale (L_hat) vs technology (t)"""
    sc = 0.5 * ((fun(L1, t0) - fun(L0, t0)) + (fun(L1, t1) - fun(L0, t1)))
    te = 0.5 * ((fun(L0, t1) - fun(L0, t0)) + (fun(L1, t1) - fun(L1, t0)))
    return sc, te
\end{lstlisting}

\subsection{AI 使用披露}

依据赛题“AI 使用规范”，现披露如下。

\textbf{所用工具：}Claude（Anthropic）、Kimi Code（Kimi）。

\textbf{使用环节与贡献范围：}(1) 资料整理与概念理解：梳理赛题附件的字段结构与《数据说明》中的条目；(2) 代码编写与调试：数据读取、向量化求解器、绘图脚本的编写与报错排查；(3) 结果核对：编写自动核对脚本，把报告中的数字与程序输出逐项比对；(4) 排版与文字：LaTeX 排版、图表整理与语句润色。上述 AI 产出的代码、数值与文字均经参赛队逐项复核。

\textbf{参赛队的工作与责任：}模型路线的选择（证据分层、质量项形式由交叉验证裁决、结构性转移按 KKT 活跃集定义、分解方法并列报告等）、公式推导与结论判断由参赛队确定，并对 AI 输出的全部代码与数值逐项运行、核验。文中理论、公式与结论均引用正式文献或由附件数据验证，未以 AI 输出作为学术依据。全部数值结果可由随文代码在附件数据上复现。
